\PassOptionsToPackage{unicode}{hyperref}
\PassOptionsToPackage{hyphens}{url}
\documentclass[11pt,a4paper]{article}

% ── Codificação e língua ──────────────────────────────────────────────────────
\usepackage[T1]{fontenc}
\usepackage[utf8]{inputenc}
\usepackage[brazil]{babel}

% ── Tipografia ────────────────────────────────────────────────────────────────
\usepackage{lmodern}
\usepackage{microtype}

% ── Cores ─────────────────────────────────────────────────────────────────────
\usepackage{xcolor}
\definecolor{javaOrange}{HTML}{E76F00}
\definecolor{javaDeep}{HTML}{BF5700}
\definecolor{javaBlue}{HTML}{1565C0}
\definecolor{javaTeal}{HTML}{00695C}
\definecolor{javaAmber}{HTML}{B45309}
\definecolor{javaInk}{HTML}{2B2140}
\definecolor{javaCodeBg}{HTML}{FFF8F0}
\definecolor{javaRule}{HTML}{FFD7B5}
\definecolor{javaShade}{HTML}{FFF3E8}

% ── Caixas e código ───────────────────────────────────────────────────────────
\usepackage{tcolorbox}
\tcbuselibrary{skins,breakable,listings}
\usepackage{listings}
\usepackage{fvextra}

\lstdefinestyle{javastyle}{
  language=Java,
  basicstyle=\ttfamily\small,
  keywordstyle=\color{javaBlue}\bfseries,
  stringstyle=\color{javaDeep},
  commentstyle=\color{javaTeal}\itshape,
  numberstyle=\tiny\color{gray},
  breaklines=true,
  breakatwhitespace=false,
  showstringspaces=false,
  tabsize=4,
}

\newtcblisting{javacode}{
  listing only,
  listing style=javastyle,
  breakable,
  enhanced,
  colback=javaCodeBg,
  colframe=javaDeep,
  leftrule=4pt,
  rightrule=0.4pt,
  toprule=0.4pt,
  bottomrule=0.4pt,
  arc=2pt,
  fontupper=\small,
  top=4pt, bottom=4pt, left=6pt, right=6pt,
}

\newtcbox{\inlinecode}{
  on line,
  colback=javaCodeBg,
  colframe=javaRule,
  boxrule=0.5pt,
  arc=2pt,
  fontupper=\ttfamily\small\color{javaDeep},
  boxsep=1pt, left=2pt, right=2pt, top=1pt, bottom=1pt,
}

% ── Layout ────────────────────────────────────────────────────────────────────
\usepackage[margin=2.4cm, top=2.8cm, bottom=2.8cm]{geometry}

% ── Cabeçalho / rodapé ────────────────────────────────────────────────────────
\usepackage{fancyhdr}
\pagestyle{fancy}
\fancyhf{}
\renewcommand{\headrulewidth}{0.4pt}
\renewcommand{\headrule}{\hbox to\headwidth{\color{javaRule}\leaders\hrule height \headrulewidth\hfill}}
\fancyhead[L]{\small\sffamily\color{javaDeep} Java Programming MOOC}
\fancyhead[R]{\small\sffamily\color{javaDeep} Resoluções · Lição 5}
\fancyfoot[C]{\small\sffamily\color{gray}\thepage}

% ── Seções ────────────────────────────────────────────────────────────────────
\usepackage{titlesec}
\titleformat{\section}
  {\sffamily\Large\bfseries\color{javaOrange}}
  {\thesection}{0.7em}{}
  [\vspace{2pt}{\color{javaRule}\titlerule[1.2pt]}]
\titleformat{\subsection}
  {\sffamily\large\bfseries\color{javaDeep}}
  {\thesubsection}{0.6em}{}
\titleformat{\subsubsection}
  {\sffamily\normalsize\bfseries\color{javaTeal}}
  {\thesubsubsection}{0.5em}{}

% ── Listas ────────────────────────────────────────────────────────────────────
\usepackage{enumitem}
\setlist[itemize]{itemsep=2pt, topsep=4pt, parsep=0pt}
\setlist[enumerate]{itemsep=2pt, topsep=4pt, parsep=0pt}

% ── Tabelas ───────────────────────────────────────────────────────────────────
\usepackage{booktabs}
\usepackage{array}
\usepackage{colortbl}
\usepackage{longtable}

% ── Hiperlinks ────────────────────────────────────────────────────────────────
\usepackage{hyperref}
\hypersetup{
  colorlinks=true,
  linkcolor=javaBlue,
  urlcolor=javaBlue,
  citecolor=javaBlue,
  pdftitle={Programação em Java --- Resoluções --- Lição 5},
  pdfauthor={Java Programming MOOC · Universidade de Helsinque},
  pdflang={pt-BR},
}

% ── Caixa de destaque (dica/atenção) ─────────────────────────────────────────
\newtcolorbox{destaque}[1][Observação]{
  enhanced, breakable,
  colback=javaShade,
  colframe=javaOrange,
  leftrule=4pt, rightrule=0.4pt, toprule=0.4pt, bottomrule=0.4pt,
  arc=2pt,
  fonttitle=\sffamily\bfseries\color{javaOrange},
  title=#1,
  top=4pt, bottom=4pt,
}

% ── Título personalizado ──────────────────────────────────────────────────────
\makeatletter
\newcommand{\subtitle}[1]{\gdef\@subtitle{#1}}
\renewcommand{\maketitle}{%
  \begin{center}
    {\color{javaRule}\rule{\linewidth}{1pt}}\\[6pt]
    {\Huge\sffamily\bfseries\color{javaOrange} \@title}\\[6pt]
    \ifx\@subtitle\undefined\else
      {\Large\sffamily\color{javaDeep} \@subtitle}\\[4pt]
    \fi
    {\small\sffamily\color{gray} \@author}\\[2pt]
    {\small\sffamily\color{gray} \@date}\\[6pt]
    {\color{javaRule}\rule{\linewidth}{1pt}}
  \end{center}
  \vspace{1em}
}
\makeatother

% ─────────────────────────────────────────────────────────────────────────────
\title{Programação em Java — Resoluções}
\subtitle{Lição 5}
\author{Java Programming MOOC · Universidade de Helsinque — tradução para o português}
\date{2026}
% ─────────────────────────────────────────────────────────────────────────────

\begin{document}
\maketitle
\tableofcontents
\vspace{1.5em}

% ─────────────────────────────────────────────────────────────────────────────
\section{Exercício 1 --- Sobrecarga de construtores}

\begin{javacode}
import java.util.Objects;

public class Livro {
    private String titulo;
    private String autor;
    private int anoPublicacao;

    public Livro(String titulo, String autor, int anoPublicacao) {
        this.titulo = titulo;
        this.autor = autor;
        this.anoPublicacao = anoPublicacao;
    }

    public Livro(String titulo, String autor) {
        this(titulo, autor, 0);
    }

    public Livro(String titulo) {
        this(titulo, "Desconhecido", 0);
    }

    @Override
    public String toString() {
        return titulo + " - " + autor + " (" + anoPublicacao + ")";
    }

    @Override
    public boolean equals(Object obj) {
        if (this == obj) return true;
        if (!(obj instanceof Livro)) return false;

        Livro outro = (Livro) obj;
        return titulo.equals(outro.titulo) && autor.equals(outro.autor);
    }

    @Override
    public int hashCode() {
        return Objects.hash(titulo, autor);
    }

    public static void main(String[] args) {
        Livro l1 = new Livro("Dom Casmurro", "Machado de Assis", 1899);
        Livro l2 = new Livro("O Cortiço", "Aluísio Azevedo");
        Livro l3 = new Livro("Sem título definido");
        Livro l4 = new Livro("Dom Casmurro", "Machado de Assis", 2020);

        System.out.println(l1);
        System.out.println(l2);
        System.out.println(l3);

        System.out.println("l1 equals l4? " + l1.equals(l4)); // true, mesmo com ano diferente
    }
}
\end{javacode}

\begin{destaque}
O uso de \inlinecode{this(...)} no início dos construtores encurtados evita repetir a lógica de atribuição --- cada construtor ``cai'' no construtor mais completo passando os valores padrão que faltam.
\end{destaque}

% ─────────────────────────────────────────────────────────────────────────────
\section{Exercício 2 --- Referências vs. primitivas}

\begin{javacode}
public class ReferenciasVsPrimitivas {

    static class Contador {
        int valor;

        Contador(int valor) {
            this.valor = valor;
        }
    }

    public static void main(String[] args) {
        // Copiando um tipo primitivo (int)
        int a = 10;
        int b = a; // b recebe uma CÓPIA do valor de a

        b = 99;

        System.out.println("--- Primitivos ---");
        System.out.println("a = " + a); // continua 10
        System.out.println("b = " + b); // 99

        // Copiando uma referência a um objeto
        Contador c1 = new Contador(10);
        Contador c2 = c1; // c2 aponta para o MESMO objeto que c1

        c2.valor = 99;

        System.out.println("--- Objetos ---");
        System.out.println("c1.valor = " + c1.valor); // também vira 99!
        System.out.println("c2.valor = " + c2.valor); // 99
    }
}
\end{javacode}

\begin{destaque}[Saída]
\begin{verbatim}
--- Primitivos ---
a = 10
b = 99
--- Objetos ---
c1.valor = 99
c2.valor = 99
\end{verbatim}
\end{destaque}

Isso mostra a diferença fundamental: \inlinecode{int} é copiado por valor, enquanto uma variável de objeto guarda uma referência (um ``endereço''), então \inlinecode{c1} e \inlinecode{c2} acabam apontando para o mesmo espaço na memória.

% ─────────────────────────────────────────────────────────────────────────────
\section{Exercício 3 --- Ponto no plano}

\begin{javacode}
public class Ponto {
    private double x;
    private double y;

    public Ponto(double x, double y) {
        this.x = x;
        this.y = y;
    }

    public double getX() {
        return x;
    }

    public double getY() {
        return y;
    }

    public double distanciaAte(Ponto outro) {
        double dx = this.x - outro.x;
        double dy = this.y - outro.y;
        return Math.sqrt(dx * dx + dy * dy);
    }

    @Override
    public boolean equals(Object obj) {
        if (this == obj) return true;
        if (!(obj instanceof Ponto)) return false;

        Ponto outro = (Ponto) obj;
        return Math.abs(this.x - outro.x) < 0.001 && Math.abs(this.y - outro.y) < 0.001;
    }

    @Override
    public String toString() {
        return "(" + x + ", " + y + ")";
    }
}
\end{javacode}

\begin{javacode}
public class Circulo {
    private Ponto centro;
    private double raio;

    public Circulo(Ponto centro, double raio) {
        this.centro = centro;
        this.raio = raio;
    }

    public boolean contemPonto(Ponto p) {
        return centro.distanciaAte(p) <= raio;
    }

    @Override
    public String toString() {
        return "Circulo{centro=" + centro + ", raio=" + raio + "}";
    }

    public static void main(String[] args) {
        Ponto centro = new Ponto(0, 0);
        Circulo circulo = new Circulo(centro, 5);

        Ponto dentro = new Ponto(3, 4); // distância = 5, está na borda
        Ponto fora = new Ponto(10, 10);

        System.out.println(circulo);
        System.out.println("(3,4) está dentro? " + circulo.contemPonto(dentro));
        System.out.println("(10,10) está dentro? " + circulo.contemPonto(fora));
    }
}
\end{javacode}

% ─────────────────────────────────────────────────────────────────────────────
\section{Exercício 4 --- Null safety}

\begin{javacode}
public class Pessoa {
    private String nome;
    private int idade;

    public Pessoa(String nome, int idade) {
        this.nome = nome;
        this.idade = idade;
    }

    public String getNome() {
        return nome;
    }

    public int getIdade() {
        return idade;
    }
}
\end{javacode}

\begin{javacode}
public class NullSafety {

    public static void imprimirPessoa(Pessoa p) {
        if (p == null) {
            System.out.println("Pessoa não encontrada.");
            return;
        }
        System.out.println(p.getNome() + ", " + p.getIdade() + " anos");
    }

    public static void main(String[] args) {
        Pessoa pessoa = new Pessoa("Maria", 30);

        imprimirPessoa(pessoa);
        imprimirPessoa(null);
    }
}
\end{javacode}

\begin{destaque}[Saída]
\begin{verbatim}
Maria, 30 anos
Pessoa não encontrada.
\end{verbatim}
\end{destaque}

\end{document}
